% =====================================================================================
% sections/s7_pitfalls.tex -- Two implementation pitfalls, demonstrated on minimal examples.
% Demonstration of Section 7.1: code/s7_pitfalls.py -> data/s7_pitfalls_autograd.json,
%   figures/s7_pitfalls.pdf; with --passes 1500 -> data/s7_pitfalls_autograd_1500passes.json.
% Demonstration of Section 7.2: research_benchmark/scripts/run_oneface.py
%   -> research_benchmark/results/oneface_laplace3d.json.
% Details, the listing and the figure: sections/S6_pitfalls.tex.  All "% src:" paths are relative
% to the root of the repository.
% =====================================================================================
\section{Two implementation pitfalls, demonstrated on minimal examples}\label{s7_pitfalls:sec}

A PINN can reach a small training loss and still be far from the solution for reasons that
have nothing to do with the optimiser or the points. This section demonstrates two such
reasons on minimal, released examples. Neither depends on seeds or budgets, and each is caught
by a cheap test that we recommend in Section~\ref{s8_discussion:sec:checklist}.

% -------------------------------------------------------------------------------------
\subsection{A residual that automatic differentiation makes identically zero}%
\label{s7_pitfalls:sec:autograd}

A tempting way to obtain $u_{tt}$ in PyTorch is to differentiate $u_t$ with respect to the
column \code{x[:, 0]} of the input tensor \code{x}. Indexing creates a new tensor, and $u_t$ was
computed from \code{x}, not from that slice, so reverse-mode differentiation~\citep{paszke2019}
finds no path between the two. Without \code{allow_unused=True} the call raises an error; with
it the call returns \code{None}, and code that replaces a missing derivative by zeros turns the
residual into the constant zero, which gives the loss neither a value nor a gradient. The correct
pattern differentiates with respect to the whole input and indexes the result
(Listing~\ref{s7_pitfalls:lst:autograd}).
% src: data/s7_pitfalls_autograd.json (autograd_demo; operator_unit_test_float64: faulty 0.0 and 0.0; correct 3.55e-15 and 19.61)
A unit test exposes the fault at once: in double precision on 2000 random points the faulty
operator for the wave equation $u_{tt}=\lap u$ returns exactly 0 both for the solution of \PWtwo\
and for the non-solution $\ee^{-t}\sin\pi x\sin\pi y$, whereas the correct operator returns at
most $\sci{3.6}{-15}$ for the former and values up to 19.6 for the latter.

% src: data/s7_pitfalls_autograd.json (training: steps 23550, none_returns 70650 of 70650 calls, max_pde_term_over_all_steps 0.0, loss 0.1011 -> 1.1007e-04; heldout_21x81x81: 1.3507 vs exact, 0.0247 vs target; slices at t = 0, 1: amplitude 0.98110, 0.36431)
Training gives no warning. We trained a 3--50--50--1 tanh network on the domain of \PWtwo\ with
a loss made of the faulty residual, a data term that prescribes $\ee^{-t}\sin\pi x\sin\pi y$ at
10\,000 random space--time points and a zero boundary term at 400 points of the spatial boundary at
$t=0$ (batches of 64, \Adam\
with learning rate $10^{-3}$, seed 0, 23\,550 iterations). The PDE term was exactly 0 at every
iteration (all 70\,650 second-derivative calls returned \code{None}), while the loss fell from
0.10 to $\sci{1.1}{-4}$. The network fitted the data term (relative distance 0.025 on the held-out
grid of \PWtwo) and followed its decay, with mode amplitudes 0.981 and 0.364 at $t=0$ and 1; its
relative distance from the solution of \PWtwo\ is 1.35. The network thus follows whatever the data
term prescribes; a correct residual would have penalised this data term, which is not a wave:

\begin{proposition}[A decaying mode is not a wave solution]\label{s7_pitfalls:prop:notwave}
\StmtNotWave
\end{proposition}

% src: data/closed_forms.json (wave2d.rel_L2_target_vs_exact_ut0_zero = 1.3800); data/s7_pitfalls_autograd.json (heldout_21x81x81.rel_L2_target_vs_exact = 1.3528); data/s7_pitfalls_autograd_1500passes.json (training: steps 235500, max_pde_term_over_all_steps 0.0, loss 0.1011 -> 5.61e-06; heldout_21x81x81: 1.3532 vs exact)
\noindent On the held-out grid, which has 21 time levels, the ratio of part~(b) is 1.353. With ten
times as many iterations nothing changes in kind: the PDE term is exactly 0 at all 235\,500
iterations, the loss falls to $\sci{5.6}{-6}$, and the distance from the solution is 1.35. Three
guards would each catch the fault: never pass \code{allow_unused=True} in a residual; assert that
every derivative is not \code{None}; and unit-test each residual operator on an exact solution
\emph{and} on a non-solution. The last is decisive, since a test on an exact solution alone is
passed by the zero operator.

% -------------------------------------------------------------------------------------
\subsection{A small loss does not imply a unique solution}\label{s7_pitfalls:sec:oneface}

A boundary-value problem with too few conditions has a whole family of solutions, and a PINN
trained on it may converge to any of them. The Laplace equation in the cube with data on one face
is the simplest example.

\begin{proposition}[A Laplace problem with data on one face only]\label{s7_pitfalls:prop:nonunique}
\StmtNonunique
\end{proposition}

% src: research_benchmark/results/oneface_laplace3d.json (oneface_laplace3d.rows: rel_l2 1.616, 12.51, 13.22, 7.275, 1.633; final_bc_loss 2.7e-04..2.2e-03, final_pde_loss 4.1e-04..7.9e-04; mean_abs_u_face_x0 0.267, 2.62, 2.65, 1.45, 0.259); research_benchmark/scripts/run_oneface.py
We trained five seeded $2\times100$ tanh networks on this problem, with the residual on the closed
$21^3$ grid of the cube and the datum on its 441 points with $x=\pi$ (AdamW~\citep{loshchilov2019}
with learning rate $10^{-3}$ and its default weight decay, 2000 iterations). Both loss terms end between $\sci{2.7}{-4}$ and $\sci{2.2}{-3}$ in
every run, and the five networks are five different, nearly harmonic functions: their relative
$L^2$ distance from the solution of \PLthree, which has the same datum and zero data on the other
five faces, is 1.62, 12.51, 13.22, 7.28 and 1.63, and the mean of $\abs{\unet}$ on the face $x=0$
is 0.27, 2.62, 2.65, 1.45 and 0.26 (Figure~\ref{s7_pitfalls:fig:demo}), in line with part~(c) of
the proposition.
% src: verify_research_benchmark/v_oneface.json (oneface|0..2: rel_l2 1.589, 12.84, 13.36); verify_research_benchmark/v_oneface.py
A re-run of seeds 0 to 2 with the separately written code \textsf{V-bench}
(Section~\ref{s3_methods:sec:verif}) gave the distances 1.59, 12.84 and 13.36. Nothing in such a
loss selects one solution, so the outcome depends on the initial weights. Count the conditions before training and say why the solution is unique; and run
more than one seed: an answer that changes with the seed while the loss stays comparably small is
the symptom.

The settings of the two demonstrations are arbitrary, were not tuned and differ from the
protocol of Section~\ref{s3_methods:sec}. The mechanism of each, a residual that is identically
zero and a loss that does not determine the solution, does not depend on them; the particular
numbers do.
